% GENERATED FILE. Edit canonical lessons, then run npm run generate:books.
% canonical-source-hash: fnv1a64:88e191ee16c86f8e
% canonical-lessons: IT-C356-curriculum, IT-C356-candidatura, IT-C356-candidarsi, IT-C356-candidato, IT-C356-assunzione

\chapter{Cerco lavoro}
\label{ch:it-cerco-lavoro}
\hlchaptermodality{356}

\begin{chapteropening}
\textbf{By the end of this chapter:} \emph{I can apply for a job: send my CV, put myself forward as a candidate and talk about being taken on.}

Complete the last lesson of chapter 356: I can apply for a job: send my CV, put myself forward as a candidate and talk about being taken on.
\end{chapteropening}

\section[il curriculum]{il curriculum — the CV, the résumé}

\label{lesson:IT-C356-curriculum}



\begin{quote}
Before the new one: say the Italian for to soften, then the Italian for to make up for.
\end{quote}

\subsection*{You'll want to know: il curriculum}
\textbf{il curriculum} — \textquotedblleft{}the CV, the résumé\textquotedblright{}.

\begin{quote}
Ho mandato il mio curriculum a dieci aziende.
\end{quote}

\begin{quote}
\emph{I sent my CV to ten companies.}
\end{quote}

Short for \emph{curriculum vitae}; the plural is the same: \emph{i curriculum}.

\begin{cousinweb}
\emph{curriculum} is taken whole from Latin \emph{curriculum} \textquotedblleft{}a running, a race, a racetrack\textquotedblright{} (Lewis \& Short), built on \emph{currere} \textquotedblleft{}to run\textquotedblright{}; \emph{curriculum vitae} is \textquotedblleft{}the path of a life\textquotedblright{}.
\end{cousinweb}

\begin{grammarlens}[title={What you've built}]
One more word for telling, explaining and describing at B1.
\end{grammarlens}

\subsection*{Your turn}
\begin{itemize}
\raggedright
  \item \emph{Say it:} \emph{il curriculum}
  \item \emph{Say it:} \emph{il curriculum}, once more
  \item \emph{Say it:} a sentence of your own with \emph{il curriculum}
  \item \emph{Recall:} say the Italian for to soften, then the Italian for to make up for, then say \emph{il curriculum} again
\end{itemize}

\begin{tcolorbox}[breakable,colback=teal!4,colframe=teal!35!black,title={Before you move on}]
What does il curriculum mean? (\textquotedblleft{}The CV, the résumé\textquotedblright{}.) Say it once more.
\end{tcolorbox}
\section[la candidatura]{la candidatura — the application (for a job); the candidacy}

\label{lesson:IT-C356-candidatura}



\begin{quote}
Before the new one: say the Italian for to make up for, then the Italian for the CV.
\end{quote}

\subsection*{You'll want to know: la candidatura}
\textbf{la candidatura} — \textquotedblleft{}the application (for a job); the candidacy\textquotedblright{}.

\begin{quote}
La candidatura va inviata entro venerdì.
\end{quote}

\begin{quote}
\emph{The application must be sent by Friday.}
\end{quote}

\emph{Presentare la propria candidatura} — to apply. \emph{Va inviata} means must be sent.

\begin{cousinweb}
\emph{candidatura} is built on \emph{candidato}, a learned borrowing from Latin \emph{candidātus} \textquotedblleft{}dressed in white\textquotedblright{}, from \emph{candidus} \textquotedblleft{}shining white\textquotedblright{} (Lewis \& Short).
\end{cousinweb}

\begin{grammarlens}[title={What you've built}]
One more word for telling, explaining and describing at B1.
\end{grammarlens}

\subsection*{Your turn}
\begin{itemize}
\raggedright
  \item \emph{Say it:} \emph{la candidatura}
  \item \emph{Say it:} \emph{la candidatura}, once more
  \item \emph{Say it:} a sentence of your own with \emph{la candidatura}
  \item \emph{Recall:} say the Italian for to make up for, then the Italian for the CV, then say \emph{la candidatura} again
\end{itemize}

\begin{tcolorbox}[breakable,colback=teal!4,colframe=teal!35!black,title={Before you move on}]
What does la candidatura mean? (\textquotedblleft{}The application (for a job); the candidacy\textquotedblright{}.) Say it once more.
\end{tcolorbox}
\section[candidarsi]{candidarsi — to apply, to put yourself forward}

\label{lesson:IT-C356-candidarsi}



\begin{quote}
Before the new one: say the Italian for the CV, then the Italian for the application.
\end{quote}

\subsection*{You'll want to know: candidarsi}
\textbf{candidarsi} — \textquotedblleft{}to apply, to put yourself forward\textquotedblright{}.

\begin{quote}
Mi sono candidata per un posto in banca.
\end{quote}

\begin{quote}
\emph{I applied for a job at a bank.}
\end{quote}

Reflexive, with \emph{essere}: \emph{mi sono candidata}. Also to stand in an election.

\begin{cousinweb}
\emph{candidarsi} is built on \emph{candidato}, a learned borrowing from Latin \emph{candidātus} \textquotedblleft{}dressed in white\textquotedblright{} (Lewis \& Short): a Roman seeking office wore a whitened toga.
\end{cousinweb}

\begin{grammarlens}[title={What you've built}]
One more word for telling, explaining and describing at B1.
\end{grammarlens}

\subsection*{Your turn}
\begin{itemize}
\raggedright
  \item \emph{Say it:} \emph{candidarsi}
  \item \emph{Say it:} \emph{candidarsi}, once more
  \item \emph{Say it:} a sentence of your own with \emph{candidarsi}
  \item \emph{Recall:} say the Italian for the CV, then the Italian for the application, then say \emph{candidarsi} again
\end{itemize}

\begin{tcolorbox}[breakable,colback=teal!4,colframe=teal!35!black,title={Before you move on}]
What does candidarsi mean? (\textquotedblleft{}To apply, to put yourself forward\textquotedblright{}.) Say it once more.
\end{tcolorbox}
\section[il candidato]{il candidato — the candidate, the applicant}

\label{lesson:IT-C356-candidato}



\begin{quote}
Before the new one: say the Italian for the application, then the Italian for to apply.
\end{quote}

\subsection*{You'll want to know: il candidato}
\textbf{il candidato} — \textquotedblleft{}the candidate, the applicant\textquotedblright{}.

\begin{quote}
Ci sono venti candidati per un solo posto.
\end{quote}

\begin{quote}
\emph{There are twenty candidates for a single job.}
\end{quote}

A woman is \emph{la candidata}.

\begin{cousinweb}
\emph{candidato} is a learned borrowing from Latin \emph{candidātus} \textquotedblleft{}dressed in white\textquotedblright{} (Lewis \& Short), from \emph{candidus} \textquotedblleft{}shining white\textquotedblright{}, because people seeking office in Rome wore a white toga.
\end{cousinweb}

\begin{grammarlens}[title={What you've built}]
One more word for telling, explaining and describing at B1.
\end{grammarlens}

\subsection*{Your turn}
\begin{itemize}
\raggedright
  \item \emph{Say it:} \emph{il candidato}
  \item \emph{Say it:} \emph{il candidato}, once more
  \item \emph{Say it:} a sentence of your own with \emph{il candidato}
  \item \emph{Recall:} say the Italian for the application, then the Italian for to apply, then say \emph{il candidato} again
\end{itemize}

\begin{tcolorbox}[breakable,colback=teal!4,colframe=teal!35!black,title={Before you move on}]
What does il candidato mean? (\textquotedblleft{}The candidate, the applicant\textquotedblright{}.) Say it once more.
\end{tcolorbox}
\section[l'assunzione]{l'assunzione — the hiring, the taking on (of staff)}

\label{lesson:IT-C356-assunzione}



\begin{quote}
Before the new one: say the Italian for to apply, then the Italian for the candidate.
\end{quote}

\subsection*{You'll want to know: l'assunzione}
\textbf{l'assunzione} — \textquotedblleft{}the hiring, the taking on (of staff)\textquotedblright{}.

\begin{quote}
L'azienda ha fermato le assunzioni per un anno.
\end{quote}

\begin{quote}
\emph{The company stopped hiring for a year.}
\end{quote}

Built on \emph{assumere} \textquotedblleft{}to hire\textquotedblright{}. \emph{Un'assunzione a tempo indeterminato} — a permanent job.

\begin{cousinweb}
\emph{assunzione} comes from Latin \emph{assūmptiō} \textquotedblleft{}a taking up, an adoption\textquotedblright{} (EtymDB), from \emph{assūmere} \textquotedblleft{}to take to oneself\textquotedblright{}, built on \emph{sūmere} \textquotedblleft{}to take\textquotedblright{}.
\end{cousinweb}

\begin{grammarlens}[title={What you've built}]
One more word for telling, explaining and describing at B1.
\end{grammarlens}

\subsection*{Your turn}
\begin{itemize}
\raggedright
  \item \emph{Say it:} \emph{l'assunzione}
  \item \emph{Say it:} \emph{l'assunzione}, once more
  \item \emph{Say it:} a sentence of your own with \emph{l'assunzione}
  \item \emph{Recall:} say the Italian for to apply, then the Italian for the candidate, then say \emph{l'assunzione} again
\end{itemize}

\begin{tcolorbox}[breakable,colback=teal!4,colframe=teal!35!black,title={Before you move on}]
What does l'assunzione mean? (\textquotedblleft{}The hiring, the taking on (of staff)\textquotedblright{}.) Say it once more.
\end{tcolorbox}
